\begin{center}
	\Huge
	Determinanter og vektorer i Maple
\end{center}


\section*{Tværvektor}
\stepcounter{section}

Ønsker vi at bestemme en vektor, der står vinkelret på en vektor, kan vi konstruere den vektor, vi kalder for tværvektoren.

\begin{defn}[Tværvektor]
	For en vektor $\vv{v}$ med koordinaterne
	\begin{align*}
		\vv{v} = 
		\begin{pmatrix}
			v_1 \\ v_2
		\end{pmatrix}
	\end{align*}
	defineret \textit{tværvektoren} til $\vv{v}$ som
	\begin{align*}
		\hat{\vv{v}} = 
		\begin{pmatrix}
			-v_2 \\ v_1
		\end{pmatrix}.
	\end{align*}
\end{defn}	

På Figur \ref{fig:tværvektor} kan vi se, tværvektoren for en vektor.

\begin{figure}[H]
	\centering
	\begin{tikzpicture}
		\begin{axis}[
			axis lines = center, 
			xmin = -1, xmax = 6.5, 
			ymin = -1, ymax = 6.5, 
			x = 1.5cm, y = 1.5cm,
			ticks = none, 
			xlabel = $x$, ylabel = $y$
		]
			\draw[-{Stealth[scale=1.5]}, thick, color = teal] (axis cs: 4,1) -- (axis cs: 6,4) node[midway, xshift = 7pt, yshift = -7pt] {$\vv{v}$};
			\draw[-{Stealth[scale=1.5]}, thick, color = teal] (axis cs: 4,1) -- (axis cs: 1,3) node[midway,xshift = 7pt, yshift = 7pt] {$\hat{\vv{v}}$};

		\end{axis}
	\end{tikzpicture}
	\caption{En vektor og dens tværvektor}
	\label{fig:tværvektor}
\end{figure}

\begin{exa}
	Vektoren
	\begin{align*}
		\begin{pmatrix}
			4 \\ 2
		\end{pmatrix}
	\end{align*}
	har vektoren 
	\begin{align*}
		\begin{pmatrix}
			-2 \\ 4
		\end{pmatrix}
	\end{align*}
	som tværvektor. 
\end{exa}

\section*{Determinanter}
\stepcounter{section}

\begin{defn}[Determinanter]
	For to vektorer $\vv{u}$ og $\vv{v}$ givet ved
	\begin{align*}
		&\vv{u} =
		\begin{pmatrix}
			u_1 \\ u_2
		\end{pmatrix}& &\textnormal{og} &
		&\vv{v} = 
		\begin{pmatrix}
			v_1 \\ v_2
		\end{pmatrix}
	\end{align*}
	 defineres \textit{determinanten} mellem $\vv{u}$ og $\vv{v}$ som
	\begin{align*}
		\det(\vv{u},\vv{v}) = u_1v_2 - u_2v_1.
	\end{align*}
	Dette skrives også til tider
	\begin{align*}
		\det(\vv{u},\vv{v}) = 
		\left|
		\begin{array}{cc}
			u_1 & v_1 \\
			u_2 & v_2
		\end{array}
		\right|
	\end{align*}
\end{defn}

\begin{exa}
	For vektorerne
	\begin{align*}
		&\vv{u} =
		\begin{pmatrix}
			2 \\ 1
		\end{pmatrix}& &\textnormal{og} &
		&\vv{v} = 
		\begin{pmatrix}
			-3 \\ 5
		\end{pmatrix}
	\end{align*}	 
	er determinanten givet ved
	\begin{align*}
		\det(\vv{u},\vv{v}) = 2\cdot 5 - 1\cdot(-3) = 10+3 = 13
	\end{align*}
\end{exa}

\begin{setn}[Areal og determinant]
	For to vektorer $\vv{u}$ og $\vv{v}$ gælder der, at arealet af deres 
	udspændende parallelogram er lig $|\det(\vv{u},\vv{v})|$.
\end{setn}

På Figur \ref{fig:determinant} kan vi se arealet, som determinanten tilsvarer

\begin{figure}[H]
	\centering
	\begin{tikzpicture}
		\begin{axis}[
			axis lines = center, 
			xmin = -1, xmax = 5, 
			ymin = -1, ymax = 5, 
			x = 1.5cm, y = 1.5cm,
			ticks = none, 
			xlabel = $x$, ylabel = $y$
		]
			\draw[-{Stealth[scale=1.5]}, thick, color = teal] (axis cs: 0,0) -- (axis cs: 2,1) node[midway, xshift = 7pt, yshift = -7pt] {$\vv{u}$};
			\draw[-{Stealth[scale=1.5]}, thick, color = teal] (axis cs: 0,0) -- (axis cs: 1,3) node[midway,xshift = 7pt, yshift = -7pt] {$\vv{v}$};
			\fill[olive,nearly transparent] (axis cs: 0,0) -- (axis cs: 2,1) -- (axis cs:3,4) -- (axis cs: 1,3) -- cycle;
			\node[color = teal] at (axis cs:1.5,2) {$|\det(\vv{u},\vv{v})|$};

		\end{axis}
	\end{tikzpicture}
	\caption{Determinant som areal}
	\label{fig:determinant}
\end{figure}

Skal vi bestemme determinanten i Maple, så skal vi skrive
\begin{align*}
	&\texttt{with(Gym):}\\
	&\texttt{$\vv{\texttt{u}}$ := <u$_{1}$,u$_{2}$>}\\
	&\texttt{$\vv{\texttt{v}}$ := <v$_{1}$,v$_{2}$>}\\
	&\texttt{det($\vv{\texttt{u}}$,$\vv{\texttt{v}}$)}
\end{align*}


\section*{Vinkel og prikprodukt i Maple}
\stepcounter{section}

Vi lægger ud med et eksempel, der viser hvordan vi definerer en vektor i Maple samt hvordan vi laver prikprodukt og bestemmer vinkler mellem vektorer.

\begin{exa}
	Vi får givet vektorerne $\vv{u}$ og $\vv{v}$ givet ved
	\begin{align*}
		&\vv{u} = 
		\begin{pmatrix} 
			2 \\ 5
		\end{pmatrix} & &\textnormal{og} & & 
		\vv{v} = 
		\begin{pmatrix}
			-4 \\ 3
		\end{pmatrix}. 
	\end{align*}
	Skal vi definere vektorerne i Maple, så skal vi skrive
	\begin{align*}
		&\texttt{with(Gym):} \\
		&\texttt{$\vv{\texttt{u}}$ := <2 , 5>} \\
		&\texttt{$\vv{\texttt{v}}$ := <-4 , 3>}
	\end{align*}
	Ønsker vi at lægge to vektorer sammen, skriver vi
	\begin{align*}
		\vv{\texttt{u}} + \vv{\texttt{v}}
	\end{align*}
	Ønsker vi at trække dem fra hinanden, skriver vi
	\begin{align*}
		\vv{\texttt{u}} - \vv{\texttt{v}}
	\end{align*}
	Og ønsker vi at skalere dem med et tal $k$, skriver vi
	\begin{align*}
		\texttt{k$\vv{\texttt{u}}$}
	\end{align*}
	Skal vi bestemme prikproduktet mellem de to vektorer, så skal vi skrive
	\begin{align*}
		\texttt{dotP($\vv{\texttt{u}},\vv{\texttt{v}}$)} 
	\end{align*}
	Ønsker vi at bestemme vinklen mellem vektorerne, så skal vi skrive
	\begin{align*}
		\texttt{vinkel($\vv{\texttt{u}},\vv{\texttt{v}}$)}
	\end{align*}
	Ønsker vi til slut at bestemme længden af en vektor, skal vi skrive
	\begin{align*}
		\texttt{len($\vv{\texttt{u}}$)}
	\end{align*}
\end{exa}

Når Maple bestemmer vinklen mellem to vektorer, så bruger den følgende resultat.
\newpage
\begin{setn}[Vinkel mellem vektorer]
	For to vektorer $\vv{u}$ og $\vv{v}$ gælder der, at 
	\begin{align*}
		\frac{\vv{u}\cdot \vv{v}}{|\vv{u}||\vv{v}|} = \cos(\theta),
	\end{align*}
	hvor $\theta$ er vinklen mellem vektorerne $\vv{u}$ og $\vv{v}$.
\end{setn}
Vi behøver dog ikke bruge denne sætning, da vi blot kan bruge \texttt{vinkel}-kommandoen i Maple.



\subsection*{Opgave 1}
Fire vektorer er givet ved
\begin{align*}
	&\vv{a} = 
	\begin{pmatrix}
		4 \\ -3
	\end{pmatrix} 
	&
	&\vv{b} = 
	\begin{pmatrix}
		-2 \\ 5
	\end{pmatrix}
	\\
	&\vv{c} = 
	\begin{pmatrix}
		6 \\ 10
	\end{pmatrix}
	&
	&\vv{d} = 
	\begin{pmatrix}
		-11 \\ -7
	\end{pmatrix}
\end{align*}

\begin{enumerate}[label=\roman*)]
	\item Bestem $4\vv{a}$.
	\item Bestem $\vv{a} + \vv{b}$.
	\item Bestem $2\vv{b} - \vv{c}$.
	\item Bestem $\vv{a}+\vv{b}+\vv{c}+\vv{d}$.
	\item Bestem $-2\vv{d} - \vv{a}$.
	\item Bestem $\vv{a}-\vv{b}+\vv{c}-3\vv{d}$.
\end{enumerate}

\subsection*{Opgave 2}

Fire punkter er givet ved $A(1,3)$, $B(-5,6)$, $C(-4,1)$ og $D(12,-7)$.
\begin{enumerate}[label = \roman*)]
	\item Bestem koordinaterne til vektorerne $\vv{AB}$, $\vv{BC}$ og $\vv{DA}$.
	\item Bestem længden af vektorerne $\vv{AB}$, $\vv{BC}$ og $\vv{DA}$
	\item Bestem $\vv{BA} + \vv{CD}$.
	\item Bestem $|2\vv{CA} - 3\vv{BC} + \vv{DB}|$.
\end{enumerate}

\subsection*{Opgave 3}

\begin{enumerate}[label=\roman*)]
\item Bestem prikproduktet og vinklen mellem følgende vektorer
\begin{align*}
&1) \ \begin{pmatrix}1 \\ 0\end{pmatrix} \textnormal{ og } \begin{pmatrix}0 \\ 1\end{pmatrix}    &&2) \  \begin{pmatrix} 4 \\  5\end{pmatrix} \textnormal{ og } \begin{pmatrix} -2  \\ 6\end{pmatrix}    \\
&3) \ \begin{pmatrix}12 \\ 15\end{pmatrix} \textnormal{ og } \begin{pmatrix}-3 \\ 14 \end{pmatrix}   &&4) \ \begin{pmatrix} -2 \\ -5\end{pmatrix} \textnormal{ og } \begin{pmatrix} 0.5 \\ -12 \end{pmatrix}     \\
\end{align*}
\end{enumerate}

\subsection*{Opgave 4}

Afgør hvilke af følgende par af vektorer, der er orthogonale

\begin{center}
\resizebox{0.45\textwidth}{!}{
\begin{tikzpicture}
	\begin{axis}[
		axis lines = center, 
		xmin = -4.5, xmax = 2.5, 
		ymin = -4.5, ymax = 2.5,
		grid,
		x = 1cm, y = 1cm,
		xtick = {-4,-3,...,2,3}, ytick = {-4,-3,...,2,3},
		xlabel = $x$, ylabel = $y$	
	]
		\draw[-{Stealth[scale=1.5]}, thick, color = teal] (axis cs: -3,-2) -- (axis cs: 2,-4);
		\draw[-{Stealth[scale=1.5]}, thick, color = teal] (axis cs: -2,2) -- (axis cs: -4,-4);
	\end{axis}
\end{tikzpicture}
}
\resizebox{0.45\textwidth}{!}{
\begin{tikzpicture}
	\begin{axis}[
		axis lines = center, 
		xmin = -1.5, xmax = 6.5, 
		ymin = -1.5, ymax = 6.5,
		grid,
		x = 1cm, y = 1cm,
		xtick = {-1,0,...,5,6}, ytick = {-1,0,...,5,6},
		xlabel = $x$, ylabel = $y$	
	]
		\draw[-{Stealth[scale=1.5]}, thick, color = teal] (axis cs: 4,4) -- (axis cs: 6,6);
		\draw[-{Stealth[scale=1.5]}, thick, color = teal] (axis cs: 5,0) -- (axis cs: 0,5);
	\end{axis}
\end{tikzpicture}
}

\end{center}

\subsection*{Opgave 5}

Vi får givet følgende to vektorer
\begin{align*}
	&\vv{u} =
	\begin{pmatrix}
		t-1 \\ 2
	\end{pmatrix} & &\textnormal{og} &
	&\vv{v} = 
	\begin{pmatrix}
		3 \\ -t + 3
	\end{pmatrix}
\end{align*}
\begin{enumerate}[label=\roman*)]
	\item Bestem prikproduktet mellem $\vv{u}$ og $\vv{v}$, hvis $t = 4$.
	\item Bestem $t$ så $\vv{u}$ og $\vv{v}$ er orthogonale. 
\end{enumerate}

\subsection*{Opgave 6}


To vektorer er givet ved
\begin{align*}
	&\vv{a} = 
	\begin{pmatrix}
		x^2 - 2 \\
		1	
	\end{pmatrix}& 
	&\textnormal{ og } &
	&\vv{b} = 
	\begin{pmatrix}
		1 \\
		x
	\end{pmatrix}
\end{align*}
\begin{enumerate}[label = \roman*)]
	\item Bestem $x$, så $\vv{a}$ og $\vv{b}$ er orthogonale.
\end{enumerate}

\newpage
\subsection*{Opgave 7}

Bestem tværvektoren for følgende vektorer 
\begin{align*}
	&a) \ 
	\begin{pmatrix}
		2 \\ 5
	\end{pmatrix}		
	&
    &b) \ 
    \begin{pmatrix}
    		-7 \\ 11
    \end{pmatrix}   \\
    &c) \
    \begin{pmatrix}
    		3 \\ -0.5
    \end{pmatrix}
    &
    &d) \
    \begin{pmatrix}
    		-\sqrt{3} \\ -27
    \end{pmatrix}
\end{align*}

\subsection*{Opgave 8}

Bestem determinanten for følgende par af vektorer

\begin{align*}
	&a) \	
	\begin{pmatrix}
		3 \\ -2
	\end{pmatrix}
	\textnormal{ og }
	\begin{pmatrix}
		5 \\ 1
	\end{pmatrix}
	&
	&b) \ 
	\begin{pmatrix}
		6 \\ 10
	\end{pmatrix}
	\textnormal{ og }
	\begin{pmatrix}
		4 \\ -7
	\end{pmatrix}
\end{align*}

\subsection*{Opgave 9}
Udregn følgende
\begin{align*}
	&a) \ 
	\left|
	\begin{array}{cc}
		5 & 7 \\
		-2 & -5
	\end{array}
	\right|
	&
	&b) \
	\left|
	\begin{array}{cc}
		-9 & 4 \\
		4 & -6
	\end{array}
	\right|
\end{align*}

\newpage
\subsection*{Opgave 10}

Beregn arealet af det skraverede område på følgende figur.

\begin{center}
	\begin{tikzpicture}
		\begin{axis}[
			axis lines = center, 
			xmin = -3, xmax = 7, 
			ymin = -3, ymax = 7, 
			x = 1.5cm, y = 1.5cm,
			grid,
			xlabel = $x$, ylabel = $y$
		]
			\draw[-{Stealth[scale=1.5]}, thick, color = teal] (axis cs: -2,-1) -- (axis cs: 2,-2) node[midway, xshift = -7pt, yshift = -7pt] {$\vv{u}$};
			\draw[-{Stealth[scale=1.5]}, thick, color = teal] (axis cs: -2,-1) -- (axis cs: -1,6) node[midway,xshift = -7pt, yshift = 7pt] {$\vv{v}$};
			\fill[olive,nearly transparent] (axis cs: -2,-1) -- (axis cs: 2,-2) -- (axis cs:3,5) -- (axis cs: -1,6) -- cycle;
		\end{axis}
	\end{tikzpicture}
\end{center}

\subsection*{Opgave 11}
Vektorerne $\vv{a}$ og $\vv{b}$ givet ved
\begin{align*}
	&\vv{a} =	
	\begin{pmatrix}
		2 \\ -4
	\end{pmatrix}
	&
	&\textnormal{og}
	&
	&\vv{b} =
	\begin{pmatrix}
		4 \\ -8
	\end{pmatrix}
\end{align*}
er parallelle, da $2\vv{a} = \vv{b}$.

\begin{enumerate}[label=\roman*)]
	\item Bestem $\det(\vv{a},\vv{b})$. Hvorfor giver resultatet god mening? Tænk på den geometriske fortolkning af determinanten
	\item Hvordan kan vi undersøge, om to vektorer er parallelle?
\end{enumerate} 


\subsection*{Opgave 12}

To vektorer $\vv{u}$ og $\vv{v}$ er givet ved
\begin{align*}
		&\vv{u} =
		\begin{pmatrix}
			t+2 \\ 5
		\end{pmatrix}
		& 
		&\textnormal{og} 
		&
		&\vv{v} =
		\begin{pmatrix}
			-t+3 \\ 4
		\end{pmatrix}
\end{align*}	
\begin{enumerate}[label=\roman*)]
	\item Bestem $\det(\vv{u},\vv{v})$, hvis $t = 6$
	\item Bestem $t$, så arealet af vektorernes udspændende parallelogram er 20.
\end{enumerate}

\subsection*{Opgave 13}
\begin{enumerate}[label=\roman*)]
	\item Bevis, at en vektor $\vv{v}$ og dens tværvektor $\hat{\vv{v}}$ er orthogonale
	\item Bevis, at der for to vektorer $\vv{u}$ og $\vv{v}$ givet ved
	\begin{align*}
		&\vv{u} =
		\begin{pmatrix}
			u_1 \\ u_2
		\end{pmatrix}
		& 
		&\textnormal{og} 
		&
		&\vv{v} =
		\begin{pmatrix}
			v_1 \\ v_2
		\end{pmatrix}
	\end{align*}
	gælder, at
	\begin{align*}
		\det(\vv{u},\vv{v}) = \hat{\vv{u}}\cdot \vv{v}.
	\end{align*}
\end{enumerate}

